\documentclass[12pt]{scrartcl}
\usepackage[a4paper,margin=3.0cm]{geometry}
\usepackage{../common/rlstudy}
\linespread{1.15}\selectfont

\title{Extended reading --- Week 01: MDPs and Bandits}
\date{\today}
\begin{document}
\maketitle

\section*{How to use this list}
Read in the order given. The first paper establishes the modern statistical
framing of the bandit problem; the second is the algorithm you will implement;
the third is the modern Bayesian counterpoint; the fourth is the contextual
generalisation that powers production matching systems; the fifth is the
unifying analysis that makes all of the above feel like one theory.

\section{Robbins (1952) --- Some aspects of the sequential design of
experiments}
\textit{Bull.\ Amer.\ Math.\ Soc.\ 58 (1952), 527--535.}

\paragraph{Summary.} Robbins frames the two-armed bandit as a problem of
\emph{minimising expected loss relative to an oracle}. He proves that a simple
``play the winner'' strategy is asymptotically consistent and outlines the
trade-off between exploration and exploitation that has structured the entire
field since. The paper predates Lai \& Robbins by 33 years and is short enough
to read in a sitting.

\paragraph{Why it matters for this week.} Robbins is the historical anchor for
the regret framework you proved on UCB1. Reading it makes the modern
$O(\log T)$ guarantees feel less like notation games and more like a slow
accumulation of the right definitions.

\begin{quote}\itshape
``We consider the question of how, on the basis of accumulating experience, to
choose between two methods of treatment so as to maximize the expected number
of `successes' or, equivalently, to minimize the expected loss\ldots\ The
problem so formulated is statistical in character because the underlying
parameters are unknown.''
\end{quote}
The two-line setup is striking: ``maximise successes'' immediately becomes
``minimise loss vs.\ the unknown best'', i.e.\ regret. The entire 1952 paper
is built on this dual.

\section{Auer, Cesa-Bianchi \& Fischer (2002) --- Finite-time analysis of the
multi-armed bandit problem}
\textit{Machine Learning 47, 235--256.}

\paragraph{Summary.} Introduces UCB1, UCB1-Tuned, and UCB2, and proves
\emph{finite-time} regret bounds of the form $O((K \log T)/\Delta)$. Before
this paper most bandit guarantees were asymptotic; the finite-time bound made
bandits a practical tool. The proof is a careful Hoeffding + union-bound
argument that you reproduced in the notes.

\paragraph{Why it matters for this week.} This is the paper whose bound you
plotted in Figure 3.2 of the notes. The proof technique --- decompose regret
into expected suboptimal pulls, bound the deviation events, sum the geometric
series --- recurs across every UCB-style algorithm in this curriculum.

\begin{quote}\itshape
``We show that, when the rewards are bounded, simple confidence-bound based
policies achieve, uniformly over time, logarithmic regret. \ldots The constants
are explicit and small enough that the algorithms can be considered for
practical purposes.''
\end{quote}
``Explicit and small enough'' is the contribution. The asymptotic theory of
Lai--Robbins was 17 years old; turning it into deployable code required the
finite-time bound to be tight in the gap $\Delta_a$.

\section{Agrawal \& Goyal (2012) --- Analysis of Thompson sampling for the
multi-armed bandit problem}
\textit{COLT 2012, arXiv:1111.1797.}

\paragraph{Summary.} First proof that Thompson sampling achieves
$O((K\log T)/\Delta)$ frequentist regret on Bernoulli bandits, matching UCB1
up to constants. The technique --- bounding the probability that a sub-optimal
arm's posterior sample exceeds the optimal sample --- is the template for
every subsequent Thompson analysis. Russo \& Van Roy (2014, NeurIPS) later
sharpened the constants via the information ratio.

\paragraph{Why it matters for this week.} Thompson sampling is the algorithm
you should reach for in production matching, and now you have a regret bound
that justifies the recommendation. The paper also empirically demonstrates
TS's superiority on the Yahoo!\ click-log dataset --- the same data Li et al.\
(2010) used for LinUCB.

\begin{quote}\itshape
``The conceptual simplicity of Thompson Sampling, along with its strong
empirical performance, makes it an attractive alternative\ldots\ However,
despite being proposed in 1933, theoretical understanding of this algorithm
has remained elusive.''
\end{quote}
``Theoretical understanding remained elusive'' for 79 years. The lesson:
empirical performance can survive without proofs for a long time, but proofs
are what enable derivative work (variance-reduction tricks, off-policy
correction, fairness constraints).

\section{Li, Chu, Langford \& Schapire (2010) --- A contextual-bandit approach
to personalized news article recommendation}
\textit{WWW 2010.}

\paragraph{Summary.} Introduces LinUCB and the offline replay evaluation
technique (now a precursor to modern OPE). Deployed at Yahoo!\ Front Page;
the replayer compares any new policy against logged behaviour without an A/B
test. The paper is the genealogical root of every production contextual-bandit
system today.

\paragraph{Why it matters for this week.} The synthetic insurance / ATS matchers you will
build in later weeks are direct descendants of this paper's architecture.
Read it for the engineering more than the maths: offline replay, hybrid
linear models, regularisation, online updating with Sherman--Morrison.

\begin{quote}\itshape
``Using only six weeks of randomly displayed events, our \texttt{LinUCB}
algorithm doubled the click rate over the editorial control. This rate
remained stable for 18 months in continuous A/B test\ldots''
\end{quote}
The bandit improvement was \emph{stable for 18 months}. That is the single
most important sentence in applied bandits: a properly-explored linear policy
does not need re-training as the world drifts because the per-step posterior
update is already an online learner.

\section{Russo \& Van Roy (2014) --- Learning to optimize via information
directed sampling}
\textit{NeurIPS 2014, arXiv:1403.5556.}

\paragraph{Summary.} Introduces the \emph{information ratio}
$\Psi_t = (\E_t[\Delta_t])^2 / I_t(A^\star; (A_t, R_t))$ --- expected
instantaneous regret squared, divided by mutual information about the optimal
action gained by the next pull. Bayesian regret is bounded by
$\sqrt{\bar\Psi \cdot H(A^\star) \cdot T}$ where $H(A^\star)$ is the prior
entropy of the optimum. The argument unifies UCB and TS proofs and motivates
\emph{information-directed sampling} (IDS) as an explicit information-gain
maximiser.

\paragraph{Why it matters for this week.} If you only learn one piece of
theory from Week 1, learn the information ratio. It makes precise the
intuition that ``a good exploration strategy explores in proportion to
information gain''. It also explains why TS works: TS implicitly bounds its
information ratio, so its Bayesian regret is automatically near-optimal.

\begin{quote}\itshape
``We propose a new approach to balancing exploration with exploitation: at each
period, sample an action so that the squared expected regret is small relative
to the information acquired about the optimal action.''
\end{quote}
The information-theoretic re-framing turns ``how much to explore?'' into a
\emph{quantity you can compute}. Subsequent work (Kirschner \& Krause 2018;
Hao, Lattimore \& Szepesv\'ari 2022) builds linear and Gaussian-process
information-ratio bounds.

\section*{Reading exercises}
After working through the five papers above, attempt the following.
\begin{enumerate}\setlength\itemsep{4pt}
  \item \textbf{Robbins, sequel.} Read Lai \& Robbins (1985, \emph{Adv.\ Appl.\
        Math.}) and reproduce the change-of-measure step that delivers the
        $1/\KL$ lower bound. Where in the proof does the assumption of
        \emph{consistency} of the algorithm enter? Could the bound be tightened
        for a restricted family of algorithms (say, only UCB-type)?

  \item \textbf{Auer et al., gap dependence.} Derive the regret bound for
        UCB1 when the gaps satisfy $\Delta_a = \Delta$ for all sub-optimal
        $a$. The bound becomes $O((K-1)\log T / \Delta)$; verify this against
        your empirical regret curve in the workbook. What is the implicit
        worst case as $\Delta \to 0$? (Hint: minimax-style $O(\sqrt{KT})$
        regret of MOSS.)

  \item \textbf{Agrawal--Goyal, beyond Bernoulli.} Extend Thompson sampling to
        Gaussian rewards with a Normal-Gamma prior. Write out the conjugate
        update and code it. Compare its regret with UCB1 on a 5-arm Gaussian
        bandit with unit-variance rewards and means in $\{0, 0.1, \dots, 0.4\}$.

  \item \textbf{Li et al., hybrid LinUCB.} Read \S 3.2 of the paper. Hybrid
        LinUCB shares a global coefficient $\beta$ across arms and a
        per-arm $\theta_a$. Write down the joint Bayesian linear-regression
        update for $(\beta, \{\theta_a\})$. Why does hybrid LinUCB dominate
        disjoint LinUCB when arms have few observations each?

  \item \textbf{Russo--Van Roy, information ratio for two-arm Bernoulli.}
        Compute the information ratio $\Psi$ of Thompson sampling on a
        two-arm Bernoulli bandit with prior Beta$(1,1)$ at $t=1$. Verify the
        bound $\Psi \le K/2 = 1$ that Russo--Van Roy prove for general
        $K$-armed bandits.
\end{enumerate}

\section*{Connections to later weeks}
The exact same machinery powers the following later weeks:
\begin{itemize}\setlength\itemsep{2pt}
  \item \emph{Week 11 (contextual bandits + OPE)} extends LinUCB to neural
        bandits and integrates the off-policy estimators from
        \texttt{research/bandits\_matching\_sota.md} \S 3.
  \item \emph{Week 12 (RL for ranking)} treats the slate-bandit problem as a
        contextual bandit over actions of size $k$, with position-bias
        correction.
  \item \emph{Week 17 (exploration / meta)} brings the information ratio back
        in the form of information-directed sampling and Bayesian deep
        exploration.
  \item \emph{Week 22 (CMDPs)} adds safety constraints to the OFU principle
        you learned this week, producing the Constrained Policy Optimisation
        / Lagrangian recipe.
\end{itemize}
The thread of ``optimism in the face of uncertainty'' (OFU) runs through
the curriculum from Week 1 to Week 24. Memorising the bandit version makes
the deep version readable.

\section*{Sourced from}
\verb|research/bandits_matching_sota.md| sections 1--3. For the broader
historical context, see Lattimore \& Szepesv\'ari, \emph{Bandit Algorithms}
(Cambridge UP, 2020), Chapters 7--19, and Sutton \& Barto, Chapter 2.

\end{document}
